\documentclass[12pt,a4paper]{article}

\usepackage[T2A]{fontenc}
\usepackage[utf8]{inputenc}
\usepackage[russian]{babel}
\usepackage{geometry}
\usepackage{enumitem}
\usepackage{listings}
\usepackage{xcolor}
\usepackage{hyperref}
\usepackage{longtable}
\usepackage{array}

\geometry{left=25mm,right=15mm,top=20mm,bottom=20mm}
\setlength{\parindent}{1.25cm}
\setlength{\parskip}{0pt}

\lstset{
    basicstyle=\ttfamily\small,
    frame=single,
    breaklines=true,
    showstringspaces=false,
    columns=fullflexible,
    keepspaces=true
}

\begin{document}

\begin{titlepage}
    \begin{center}
        Федеральное государственное автономное образовательное учреждение\\
        высшего образования\\[2mm]
        \textbf{«Национальный исследовательский университет\\
        ``Высшая школа экономики''»}\\[3mm]
        Факультет компьютерных наук\\
        Департамент программной инженерии

        \vfill

        {\Large \textbf{Отчёт по индивидуальному заданию № 1}}\\[4mm]
        {\large по дисциплине «Архитектура компьютера и операционные системы»}\\[8mm]
        {\Large \textbf{Вариант 14. «Задача об умных цветах»}}

        \vfill

        \begin{flushright}
            Выполнила: Гончарова Екатерина\\[3mm]
            Группа: БПИ-253
        \end{flushright}

        \vfill
        Москва, 2026
    \end{center}
\end{titlepage}

\tableofcontents
\newpage

\section{Постановка задачи}

Требуется разработать консольное приложение на языке C для ОС Linux,
которое с помощью последовательной событийной имитации моделирует
взаимодействие умных цветов и садовников.

На клумбе находится заданное количество цветов. Каждый цветок самостоятельно
изменяет своё состояние и при недостатке влаги формирует заявку на полив.
Садовники принимают активные заявки и обслуживают их. Если несколько садовников
могут откликнуться на одну заявку, заявка должна быть закреплена только за одним
садовником. После полива действует защитный интервал, в течение которого
повторный полив считается переливом и может привести к гибели цветка.

В модели используются следующие состояния цветка:
\begin{itemize}[noitemsep]
    \item нормальное состояние;
    \item увядание;
    \item полив;
    \item восстановление;
    \item засыхание;
    \item гибель от перелива.
\end{itemize}

Программа должна отображать изменение состояния цветов, создание и снятие
заявок, принятие заявки садовником, начало и окончание полива, восстановление,
засыхание и перелив. Моделирование завершается после истечения заданного
периода, после прекращения активности всех цветов или по запросу пользователя.

\section{Входные данные}

Набор входных параметров программы соответствует перечню из условия варианта:
\begin{itemize}[noitemsep]
    \item количество цветов;
    \item индивидуальные скорости высыхания;
    \item порог формирования заявки на полив;
    \item порог засыхания;
    \item длительность полива;
    \item длительность восстановления;
    \item защитный интервал от повторного полива;
    \item количество садовников;
    \item скорость работы садовников;
    \item правило выбора заявок.
\end{itemize}

Дополнительные пользовательские параметры в программу не вводятся. Начальная
влажность и период наблюдения являются внутренними константами модели. После
успешного завершения полива влажность цветка восстанавливается до максимального
значения.

Параметры могут вводиться интерактивно или считываться из файла. Для тестовых
запусков используются файлы формата \texttt{.cfg} с парами \texttt{key=value}.
Например:

\begin{lstlisting}
flower_count=6
drying_rates=4,7,5,9,6,8
request_threshold=47
dry_threshold=14
watering_duration=2
recovery_duration=2
protection_interval=5
gardener_count=2
gardener_speed=2
request_policy=first
\end{lstlisting}

Для правила выбора заявки поддерживаются два значения: \texttt{first} --- первая
доступная заявка по номеру цветка и \texttt{driest} --- заявка активного цветка
с наименьшей текущей влажностью.

\section{Модель и сценарий работы}

Один проход главного цикла программы соответствует одному условному такту.
Внутри такта действия выполняются последовательно, но состояние каждой сущности
хранится отдельно. Это позволяет имитировать независимое поведение цветов и
садовников без использования потоков.

На каждом такте выполняются следующие действия:
\begin{enumerate}[noitemsep]
    \item для каждого активного цветка уменьшается оставшийся защитный интервал;
    \item цветы, которые не находятся в состоянии полива или восстановления,
          теряют влагу согласно своей индивидуальной скорости высыхания;
    \item при достижении порога заявки цветок переходит в состояние увядания и
          создаёт заявку, если активной заявки ещё нет;
    \item при достижении порога засыхания цветок прекращает активность;
    \item свободные садовники выбирают ещё не назначенные заявки;
    \item скорость работы садовника определяет, как быстро он завершает подход к выбранному цветку;
    \item после подхода начинается полив;
    \item по завершении полива влажность восстанавливается до максимума, заявка закрывается,
          включается защитный интервал и начинается восстановление;
    \item после обработки всех участников проверяется условие завершения.
\end{enumerate}

Правило выбора заявки задаётся входным параметром. В режиме \texttt{first}
садовник выбирает первую доступную заявку по номеру цветка, а в режиме
\texttt{driest} --- заявку цветка с минимальной текущей влажностью.

\section{Представление сущностей}

Для цветка используется структура \texttt{Flower}. Она хранит идентификатор,
влажность, индивидуальную скорость высыхания, текущее состояние, наличие
активной заявки, назначенного садовника и таймеры полива, восстановления и
защитного интервала.

Для садовника используется структура \texttt{Gardener}. Она хранит идентификатор,
признак занятости, индекс обслуживаемого цветка, оставшееся время текущего
действия и скорость работы.

Общие параметры одного запуска объединены в структуру \texttt{Settings}.
Это позволяет передавать настройки между функциями одной переменной и не
использовать глобальные параметры модели.

\section{Структура программы}

Проект разбит на небольшие модули по ответственности:

\begin{longtable}{|>{\raggedright\arraybackslash}p{4cm}|p{10cm}|}
\hline
\textbf{Файл} & \textbf{Назначение} \\
\hline
\texttt{src/main.c} & Минимальная точка входа программы. \\
\texttt{src/app.c} & Запуск приложения, выбор способа ввода, проверка параметров, настройка журнала и обработчика \texttt{Ctrl+C}. \\
\hline
\texttt{src/input.c} & Интерактивный ввод и вывод параметров. \\\texttt{src/config\_parser.c} & Чтение и разбор \texttt{.cfg}-файлов. \\\texttt{src/input\_validation.c} & Проверка корректности входных параметров. \\
\hline
\texttt{include/constants.h} & Именованные константы модели: границы влажности, максимальные размеры, специальные значения. \\
\texttt{src/model.c} & Инициализация цветов и садовников, базовые операции модели. \\
\hline
\texttt{src/flower\_logic.c} & Изменение влажности и состояний цветка, создание и снятие заявок. \\
\hline
\texttt{src/gardener\_logic.c} & Выбор заявок, подход к цветку, полив и освобождение садовника. \\
\hline
\texttt{src/simulation.c} & Главный цикл последовательной имитации. \\
\hline
\texttt{src/statistics.c} & Формирование итоговой статистики. \\
\hline
\texttt{src/logger.c} & Вывод событий в консоль и сохранение журнала в файл. \\
\hline
\texttt{include/} & Заголовочные файлы модулей. \\
\hline
\texttt{tests/} & Подготовленные тестовые конфигурации. \\
\hline
\end{longtable}

Такое разбиение позволяет отделить предметную логику цветов и садовников от
ввода, журналирования и управляющего цикла.

\section{Пользовательский интерфейс}

Программа является консольной. Без аргументов командной строки параметры
последовательно запрашиваются у пользователя. После проверки введённых значений
программа печатает параметры запуска и начинает моделирование.

Во время работы выводятся сообщения вида:

\begin{verbatim}
[07] Цветок 2: NORMAL -> WILTING
[07] Цветок 2: отправлена заявка на полив
[07] Садовник 1: принял заявку цветка 2
[08] Садовник 1: начал полив цветка 2
[10] Садовник 1: закончил полив цветка 2, влажность 24 -> 79
\end{verbatim}

Тот же журнал сохраняется в файл \texttt{simulation.log}. Благодаря этому
результат можно просмотреть после завершения программы.

\section{Варианты запуска}

Сборка проекта выполняется командой:

\begin{lstlisting}
make
\end{lstlisting}

Интерактивный запуск:

\begin{lstlisting}
./smart_flowers
\end{lstlisting}

Запуск с заранее подготовленным набором параметров:

\begin{lstlisting}
./smart_flowers tests/test1.cfg
\end{lstlisting}

Также в \texttt{Makefile} предусмотрены готовые цели:

\begin{lstlisting}
make test1
make test2
make test3
\end{lstlisting}

Нормальное завершение происходит после истечения внутреннего периода наблюдения либо
если все цветы прекратили активность. Досрочное завершение выполняется по
\texttt{Ctrl+C}: обработчик сигнала устанавливает флаг остановки, после чего
главный цикл корректно завершается и печатает итоговую статистику.

\section{Инварианты модели}

В реализации контролируются основные ограничения задачи:
\begin{itemize}[noitemsep]
    \item активная заявка закрепляется не более чем за одним садовником;
    \item садовник одновременно работает не более чем с одним цветком;
    \item цветок не создаёт новую заявку, пока предыдущая заявка активна;
    \item завершённая заявка снимается и не может быть обслужена повторно;
    \item засохший или погибший от перелива цветок больше не участвует в моделировании;
    \item параметры запуска проверяются до начала моделирования.
\end{itemize}

\section{Журналирование и системный ввод-вывод}

Для журнала используется файловый дескриптор. Файл открывается системной
функцией \texttt{open}, события записываются функцией \texttt{write}, а после
завершения дескриптор закрывается функцией \texttt{close}. Один и тот же текст
события отправляется в стандартный вывод и в файл журнала.

Это позволяет сохранить наглядный консольный интерфейс и одновременно
использовать ввод-вывод, основанный на файловых дескрипторах.

\section{Тестирование}

В проект включены три тестовых набора.

\subsection*{Тест 1}
Шесть цветов и два садовника с различными скоростями высыхания. Набор проверяет
обычный режим работы: постепенное появление заявок, их распределение между
садовниками, полив и последующее восстановление цветов.

\subsection*{Тест 2}
Семь цветов и три садовника. Повышенные скорости высыхания и более высокий
порог формирования заявки приводят к появлению нескольких заявок за небольшой
промежуток времени. Тест проверяет обслуживание нескольких активных заявок.

\subsection*{Тест 3}
Пять цветов и один садовник. Невысокая скорость работы единственного садовника создаёт очередь
из заявок и позволяет проверить ожидание обслуживания, засыхание отдельных
цветов и корректное завершение моделирования.

После каждого запуска программа выводит число живых, засохших и погибших от
перелива цветов, общее количество поливов, а также итоговое состояние,
влажность и число поливов каждого цветка.

\section{Дополнительные возможности}

Помимо основной логики задачи в проекте реализованы:
\begin{itemize}[noitemsep]
    \item комментарии и Doxygen-документация для структур и функций;
    \item генерация HTML-документации командой \texttt{make docs};
    \item запись журнала выполнения в \texttt{simulation.log};
    \item корректное завершение по \texttt{Ctrl+C};
    \item подготовленные тестовые наборы и команды их запуска;
    \item чтение тестовых параметров из конфигурационных файлов;
    \item проверка корректности входных параметров;
    \item итоговая статистика по состоянию цветов.
\end{itemize}

\section{Заключение}

В результате разработана последовательная имитационная программа на языке C,
моделирующая взаимодействие умных цветов и садовников. Состояние участников
хранится независимо, а их действия выполняются по тактам в определённом порядке.
Программа отражает основные события предметной области, контролирует ключевые
ограничения, поддерживает интерактивный и файловый ввод, сохраняет журнал работы
и корректно завершает выполнение как штатно, так и по пользовательскому
прерыванию.

\end{document}
